\documentclass[12pt,a4paper]{article}

\hbadness=10000

% Paquetes

\usepackage[utf8]{inputenc}
\usepackage[table]{xcolor}
\usepackage[spanish]{babel}
\usepackage[T1]{fontenc}
\usepackage[margin=2.5cm]{geometry}
\usepackage{fancyhdr}
\usepackage{tabularx}
\usepackage{setspace}
\usepackage[spanish]{babel}
%Paquetes para simbologia%
\usepackage{enumitem}  
\usepackage{cellspace} % para espaciado automático en celdas
\setlength\cellspacetoplimit{12pt}
\setlength\cellspacebottomlimit{12pt} 
\usepackage{amsmath}
\usepackage{array}
\usepackage{amsfonts}
\usepackage{relsize}
\usepackage{amssymb}
\usepackage{physics}
\usepackage{longtable}
\usepackage{graphicx}
\usepackage{wrapfig}
\usepackage{tabularx}
\usepackage{caption}
\usepackage{float}
\usepackage{xurl}
\usepackage{booktabs}
\usepackage{tikz}
\usepackage{array}
\usepackage[normalem]{ulem}
\usepackage{amsmath, amssymb}
\usepackage{fancyhdr}
\usepackage[colorlinks=true,
            linkcolor=black,
            urlcolor=myblue,
            citecolor=black,
            filecolor=black]{hyperref}
 % Opción estándar para enlaces
\Urlmuskip=0mu plus 1mu          % Mejora el espaciado para permitir cortes
\usepackage{subcaption}  % en el preámbulo
\usepackage{pgfplots}
\pgfplotsset{compat=1.18}
\usepackage{tikz}
\usepackage{xcolor}
\definecolor{myblue}{RGB}{42, 127, 179}

\pagestyle{fancy}
\fancyhf{}
\fancyhead[L]{Ingeniería Mecánica}
\fancyhead[C]{Materiales Metálicos} 
\fancyhead[R]{UTN-FRVM}
\fancyfoot[C]{\thepage}

%---------------------------------------
\setlength{\headheight}{15pt}

\begin{document}

% Portada

\begin{titlepage}
    \centering
    {\Large \textbf{UNIVERSIDAD TECNOLÓGICA NACIONAL}}\\
    {\Large \textbf{FACULTAD REGIONAL VILLA MARÍA}}\\[2cm]
    {\LARGE \textbf{Ingeniería Mecánica - Materiales Metálicos}}\\[2cm]
    {\huge \textbf{Trabajo Práctico 6-01}}\\[1.5cm]

\raggedright
Grupo DEL RIO:
\begin{itemize}
    \item Antico, Rodrigo.
    \item Cabral, Franco.
    \item Cardozo, Martín.
    \item Córdoba, Nathan.
    \item Cucco, Ramiro.
    \item del Río, Juan.
    \item Guerini, Nazareno.
    \item Medina, Ivo.
    \item Ortiz, Gastón.
    \item Picos, Elías.
    \item Quinteros, Lautaro
\end{itemize}
 Docentes:
 \begin{itemize}
     \item Dr. Lucioni, Eldo José.
     \item  Ing. Victorio Vallaro, Juan Manuel.\\[1cm]
 \end{itemize}
 
\begin{center}
18 de agosto de 2025
\end{center}
 
    \vfill
\end{titlepage}

\newpage

{\Large{\textbf{TP 6-01: Fallas de Material}}}\\[0.1cm]

{\Large{\uline{Objetivos}}:}

\begin{itemize}
  \renewcommand{\labelitemi}{}
  \item PASO 1: Efectúe la búsqueda de las palabras clave (Tabla 1) en la 
bibliografía detallada en el listado "Bibliografía para Búsqueda".
  \item PASO 2: Analice, investigue y reflexione sobre la relación de cada 
párrafo -en el que aparece la palabra clave- con el fenómeno de Falla 
de Material.
  \item REQUERIMIENTO: Estar en capacidad de explicar la causa y el efecto 
de la Falla de Material desarrollada en dicho párrafo.
\end{itemize}

\vspace{0.1cm}

\begin{table}[H] % o [htbp]
\centering
\setlength{\tabcolsep}{1.5cm}
\renewcommand{\arraystretch}{1.3}

\begin{tabular}{|c@{\hspace{1cm}}c|}
\hline
\textbf{Español} & \textbf{Inglés} \\ \hline
falla / fallar & fail / failure \\
fallo & fail / failure \\
rompe / romper / romperse  & break \\
rotura  & rip \\
roto & broken \\
fractura & fracture \\
defecto& defect  \\
imperfección  &imperfection \\
deficiencia  & deficiency \\
fatiga & fatigue \\
desgaste & wear \\
corrosión  & corrosion \\
grieta & crack \\
agrietamiento & cracking  \\
pierde & misses \\
perder & to lose  \\
\hline
\end{tabular}

\vspace{5pt}
\textbf{Tabla 1:} Lista de palabras % ← escribís lo que quieras
\end{table}

\vspace{0.1cm}

{\Large{\uline{Bibliografía para Búsqueda}}:}

\begin{itemize}
    \item  Libro 1: Apraiz Barreiro, J. - Tratamientos térmicos de los aceros. 1983. 8va Ed 
    \item  Libro 2: Callister, D. W. - Introducción a la Ciencia e Ingeniería de los Materiales - 1995
    \item  Libro 3: Shigley, Budyna, Nisbett - Diseño en ingeniería mecánica de Shigley, 9na Edición
    \item  Libro 4: Sturla, A.E. - Tratamientos Térmicos de los Aceros - Tomo II
\end{itemize}

\newpage

\section{Tabla 2 de Resultados de Búsqueda }

\begin{center}
\setlength{\tabcolsep}{4pt}
\renewcommand{\arraystretch}{1.2}

\begin{tabular}{|*{17}{c|}}
\hline
\textbf{} 
& \rotatebox{90}{\textbf{falla / fallar }} 
& \rotatebox{90}{\textbf{fallo}} 
& \rotatebox{90}{\textbf{rompe / romper / romperse}} 
& \rotatebox{90}{\textbf{rotura}} 
& \rotatebox{90}{\textbf{roto}} 
& \rotatebox{90}{\textbf{fractura}} 
& \rotatebox{90}{\textbf{defecto}} 
& \rotatebox{90}{\textbf{imperfección}} 
& \rotatebox{90}{\textbf{deficiencia}} 
& \rotatebox{90}{\textbf{fatiga}} 
& \rotatebox{90}{\textbf{desgaste}} 
& \rotatebox{90}{\textbf{corrosión}} 
& \rotatebox{90}{\textbf{grieta}} 
& \rotatebox{90}{\textbf{pierde}}
& \rotatebox{90}{\textbf{perder}} \\ \hline

\textbf{Libro 1} & 3 & - & 27/5/4 & 72 & 2 & 28 & 6 & - & - & 45 & 13 & 6 & 20 & 6 & 0 \\ \hline
\textbf{Libro 2} & -/1 & 8 & 4/6/9 & 27 & - & 49 & 34 & 10 & - & 16 & 8 & 38 & 16 & - & - \\ \hline
\textbf{Libro 3} & 374/10 & 3 & 1/-/- & 7 & 1 & 105 & 12 & 7 & 2 & 764 & 191 & 51 & 75 & 6 & - \\ \hline
\textbf{Libro 4} & 2/- & - & 7/2/- & 74 & 1 & 10 & 12 & 1 & - & 19 & 5 & 13 & 8 & - & - \\ \hline
\end{tabular}
\end{center}

A continuación citaremos 6 o más ejemplos, aproximadamente, de cada palabra del libro, que se nos asignó como equipo este año, para mostrar en dicha sección el significado o contexto de dicha palabra, para el conteo icluiremos todas las palabras pero para los ejemplos dejaremos de lado aquellas formen parte de un índice, título, subtítulo, pregunta o referencia.

El libro asignado para el caso de este año 2025 es el de  Callister, D. W. - Introducción a la Ciencia e Ingeniería de los Materiales - 1995; hay que aclarar que también extraimos algunas palabras del libro Callister, D. W. - Introducción a la Ciencia e Ingeniería de los Materiales 2da edición (Correspondiente a la 9na edición original).\\[0.1cm]

    
\section{Libro 2: Callister, D. W. - Introducción a la Ciencia e Ingeniería de los Materiales - 1995}

\textbf{Fallo:} 

subsección 4.5.4; primer párrafo: "[...] son los fallos de apilamiento, [...]"
    
Resumen Capítulo 8; tercer párrafo: "[...] y permite diseñar estructuras con una probabilidad de fallo mínima. [...]"
    
Resumen Capítulo 8; "[...] Finalmente, la deformación se acelera en la fluencia terciaria, justo antes del fallo [...]"
    
subsección 16.9.2; primer párrafo: "Los polímeros pueden experimentar fallos por fatiga [...]"
    
subsección 16.9.2; primer párrafo: "[...] el esfuerzo para que se produzca el fallo [...]"

subsección 18.7.8; : "[...], como consecuencia puede aparecer un fallo. La apariencia del fallo es el característico [...]"

Cap 10–Pág 285; cuarto párrafo; "[...]" experimentó una descompresión explosiva y un fallo estructural el 28 de abril de 1988. Una investigación del accidente concluyó
que la causa fue la fatiga del metal, agravada por corrosión por fisuración "[...]"

Cap 10–Pág 286; primer párrafo; "[...]" El diseño de un componente estructural requiere conocimientos de ingeniería para reducir al mínimo la posibilidad de fallo "[...]"

    
subsección 18.7.4; primer párrafo: "[...] pérdida de material hasta que ocurre un fallo. [...]"\\[0.1cm]
    
\textbf{Rompe:}

sección 7.5; último párrafo: "[...] normalmente se rompe antes de deformarse plásticamente.[...]"
    
sección 16.2; segundo párrafo: "[...], ya que se rompe cuando se deforma elásticamente."[...]"
    
subsección 17.5.2; segundo párrafo: "[...] donde (TS), y (T5),, representan, respectivamente, la resistencia a la fractura de la fibra y la tensión en la matriz cuando el material compuesto se rompe[...]"
    
sección 17.11; tercer párrafo: "El compuesto se rompe al fallar la fase matriz."\\[0.1cm]


\textbf{Romper:}

sección 4.7; segundo párrafo: "[...] que logran romper gran número de enlaces atómicos.[...]"
    
sección 5.2; primer párrafo: "[...] suficiente energía como para romper los enlaces con los átomos [...]"
    
subsección 16.4.1; primer párrafo: "[...] que es capaz de romper gran número de enlaces [....]"
    
sección 16.5; primer párrafo: "[...] son tan violentas que pueden romper los enlaces covalentes. [...]"
    
subsección 16.11.3; primer párrafo: "[...] y puede romper algunos de ellos [...]"
    
sección 17.2: "[...]. De este modo para romper un miembro de hormigón [...]"\\[0.1cm]

\textbf{Romperse:}

sección 8.8; tercer párrafo: "[...] el material es capaz de aguantar antes de romperse. [...]"
    
Problema 8.10; primer párrafo: "[...] Un componente de Mg0 no debe romperse cuando se aplica una tensión de 13,5 MPa (1969 psi). [...]"
    
sección 12.2; tercer párrafo: "[...] puesto que son moderadamente dúctiles y capaces de experimentar deformación permanente sin romperse. [...]"
    
sección 16.8; tercer párrafo: "[...] ,pero se convierten en dúctiles al aumentar la temperatura y aproximarse a la temperatura de transición vítrea y presentan comportamiento plástico antes de romperse. [...]"
    
sección 16.16.2; Cuarto párrafo: "[...] Al generarse una buena unión puede romperse antes el material adherido que el adhesivo. [...]"

sección 23.6; primer párrafo: "[...] La articulación puede romperse, lo cual ocurre normalmente en la región estrecha, [...]"

sección 23.3-Ecuacion 23.11; primer párrafo: "[...] para que la viga funcione correctamente, no debe romperse cuando se aplica la tensión. [...]"\\[0.1cm]
    
\textbf{Rotura:}
 
Problema 6.2; Segundo párrafo: "[...] la deformación plástica corresponde a la rotura de los enlaces entre los átomos vecinos más próximos [...]"
    
sección 6.6.2; cuarto párrafo: "[...] La tensión de fractura o bien de rotura corresponde a la tensión en la fractura. [...]"
    
sección 6.6.3; segundo párrfo: "[...] Siempre que una parte significativa de la deformación plástica a la rotura esté confinada en la región dela estricción, [...]"
    
sección 6.11; cuarto párrafo: "[...] "¿Cuál es la probabilidad de rotura de esta aleación en determinadas circunstancias?“ [...]"
    
sección 6.12; quinto párrafo: "[...] lo más importante: las consecuencias de la rotura en términos de vidas humanas y daños materiales.[...]"

sección 8.2; cuarto párrafo: "[...] La tendencia a la rotura frágil de los materiales normalmente dúctiles se discute en la Sección 8.6.[...]"

sección 8.3-Pagina 22; segundo párrafo: "[...] esta estructura es característica de la fractura que resulta de la rotura a tracción uniaxial.[...]"

sección 8.3-Pagina 22; quinto párrafo: "[...] Esta forma parabólica puede ser indicativa de la rotura por cizalladura.[...]"

sección 8.3-Pagina 24; segundo párrafo: "[...] ,la propagación de la grieta corresponde a la sucesiva y repetida rotura de enlaces atómicos a lo largo deplanos cristalográficos;[...]"

sección 8.5.1-Pagina 27; segundo párrafo: "[...] la rotura ocurrirá cuando la resistencia cohesiva teórica del material sea superada en la punta de uno de los defectos. [...]"
 
sección 8.5.3-Pagina 32; tercer párrafo: "[...] tienen valores pequeños de-K, y son vulnerables a la rotura catastrófica.[...]"

sección 8.6.2-Pagina 38; tercer párrafo: "[...] dúctil-frágil debenutilizarse únicamente a temperaturas por encima de la temperatura de transición para evitar la rotura frágil y catastrófica.[...]"

sección 8.6.2-Pagina 40; primer párrafo: "[...] La fatiga es una forma de rotura que ocurre en estructuras sometidas a tensiones dinámicas y fluctuantes [...]" 

sección 8.6.2-Pagina 40; cuarto párrafo: "[...] La fatiga es importante ya que es la primera causa de rotura de los materiales.[...]"

sección 8.6-Pagina 41; primer párrafo: "[...] La rotura por fatiga tiene aspecto frágil aun en metales que son normalmente dúctiles,[...]"

sección 8.8; tercer párrafo: "[...] S, frente al logaritmo del número N de ciclos hasta la rotura para cada una de lás probetas.[...]"

sección 8.8; tercer párrafo: "[...] ,por debajo del cual la rotura por fatiga no ocurrirá.[...]"

sección 8.9; primer párrafo: "[...] El proceso de rotura por fatiga está caracterizado por tres etapas distintas:[...]" 

sección 8.11.2; tercer párrafo: "[...] El efecto neto es que la probabilidad de nucleación de la grieta,y por tanto de rotura por fatiga se reduce.[...]"

 sección 8.12; primer párrafo: "[...] La rotura que ocurre por la acción simultánea de una tensión cíclica y el ataque químico se denomina fatiga con corrosión [...]"
 
sección 8.15-Problema 8.3; primer párrafo: "[...], predecir el tiempó a la rotura de un componentes que es sometido a una tensión [...]"

sección 8.16-Pagina 4; prime párrafo: "[...]La fatiga es un tipo de fractura que conduce a la rotura catastrófica cuando se aplican cargas fluctuantes con el tiempo.[...]"

Cap 10–Pág.287; octavo párrafo; "[...]" en una fractura dúctil la presencia de deformación plástica es una evidenciade que la rotura es inminente, permitiendo que se tomen medidas preventivas "[...]" 

sección 19.21; segundo párrafo: "[...] la corriente a través del dieléctrico por el movimiento de estos electrones aumenta drásticamente; aveces la fusión localizada,  quemadura o bien la vaporización produce la degradación irreversible y quizás incluso la rotura del material [...]"
    
sección 8.16-Pagina 4; prime párrafo: "[...]La fatiga es un tipo de fractura que conduce a la rotura catastrófica cuando se aplican cargas fluctuantes con el tiempo.[...]"\\[0.1cm]
    
\textbf{Fractura:}

cap 4-pag 74; ¿POR QUÉ ESTUDIAR Estructura en sólidos cristalinos?; primer párrafo: "[...] que tienen
una determinada estructura cristalina, son mucho más frágiles (es decir, se fracturan con menores grados de deformación) que oro y plata, puros y no deformados [...]"

Seccion 8.2-pag 210: tercer párrafo: "[...] y también para reducir la probabilidad de fractura en los extremos de la muestra [...]"

Cap 10-285: primer párrafo: "[...] Este fenómeno está relacionado con uno delos principios básicos de la mecánica de la fractura: [...]"

Cap 10-pag 285; ¿POR QUÉ ESTUDIAR Rotura?; primer párrafo: "[...] Por tanto, es importante
comprender la mecánica de los distintos modos de rotura (fractura, fatiga y fluencia) y, además, conocerlos principios de diseño [...]"

Cap 10-pag 285; Objetivo de aprendizaje; primer pregunta: "[...] Describir el mecanismo de propagación de grietas tanto en modo de fractura dúctil como frágil. [...]"

Sección 10.2-pag 287; primer párrafo: "[...] Una fractura simple es la separación de un cuerpo en dos o más partes en respuesta a una tensión estática impuesta (es decir, constante o que varía lentamente con el tiempo) [...]"\\[0.1cm]

\textbf{Defecto:}

Cap 6-pag 143; primer párrafo: "[...] gracias a los defectos atómicos es posible reducir las emisiones de gases contaminantes [...]"

Sección 6.1-pag 144; primer párrafo: "[...] Sin embargo, el sólido ideal no existe: todos los materiales tienen un gran número de defectos o imperfecciones [...]"

Sección 6.1-pag 144; segundo párrafo: "[...] Un defecto cristalino se refiere a una irregularidad de la red para la que una o más de sus dimensiones son del orden de un diámetro atómico [...]"

Sección 6.2-pag 145; tercer párrafo: "[...] Un defecto autointersticial se forma cuando un átomo del cristal se coloca en un lugar intersticial de la red, que es un pequeño espacio vacío que ordinariamente no está ocupado. [...]"

Sección 7.2-pag 182; primer párrafo: "[...] Naturalmente, este proceso necesita de la presencia de vacantes y las posibilidades de la difusión de las vacantes es función del número de estos defectos. [...]"

Cap 10–Pág.287; octavo párrafo; "[...]" en una fractura dúctil la presencia de deformación plástica es una evidenciade que la rotura es inminente, permitiendo que se tomen medidas preventivas "[...]" 

Sección 9.1-pag 254; segundo párrafo: "[...] la resistencia mecánica se podía explicar por la existencia de un tipo de defecto lineal cristalino que, desde entonces, se ha llamado dislocación [...]"\\[0.1cm]

\textbf{Imperfección:}

Cap 6-pag 143; ¿POR QUÉ ESTUDIAR Imperfecciones en sólidos?; primer párrafo: "[...] Las propiedades de algunos materiales están profundamente influenciadas por la presencia de imperfecciones. [...]"

Sección 6.1-pag 144; primer párrafo: "[...] Sin embargo, el sólido ideal no existe: todos los materiales tienen un gran número de defectos o imperfecciones [...]"

Sección 6.1-pag 144; segundo párrafo: "[...] La clasificación de las imperfecciones cristalinas normalmente se hace según la geometría o las dimensiones del defecto. [...]"\\[0.1cm]

\textbf{Fatiga:}

Cap 7-pag 180; primer párrafo: "[...] Además, dentro de esta región se introducen esfuerzos de compresión residuales, que dan lugar a una mejora en la resistencia a la rotura por fatiga del engranaje cuando está en servicio [...]"

Cap 7-pag 181; ¿POR QUÉ ESTUDIAR Difusión?; primer párrafo: "[...] es decir, su dureza y su resistencia a la rotura por fatiga se han incrementado por difusión de un exceso de carbono o nitrógeno en la capa exterior de la superficie. [...]"

Cap 10-pag 285; tercer párrafo: "[...] Una investigación del accidente concluyó
que la causa fue la fatiga del metal, agravada por corrosión por fisuración [...]"

Sección 10.1-pag 286; segundo párrafo: "[...] En este capítulo se tratan los siguientes temas: fractura simple (tanto en modo dúctil como frágil), fundamentos de mecánica de la fractura, ensayos de tenacidad de fractura, transición dúctil-frágil, fatiga y termofluencia. [...]"

Cap 10–Pág 285; cuarto párrafo; "[...]" Un programa de mantenimiento, debidamente ejecutado por la compañía aérea habría detectado el daño por fatiga y se hubiera prevenido este accidente "[...]"

Sección 10.2-pag 287; primer párrafo: "[...] La fractura también puede ocurrir por fatiga (cuando se imponen esfuerzos cíclicos) [...]"

Cap 10-pag 286; Objetivos de aprendizaje: "[...]"6. Definir el concepto de fatiga y especificar las
condiciones en que se produce. [...]"\\[0.1cm]

\textbf{Desgaste:}

Cap 7-pag 180; primer párrafo: "[...] El aumento del contenido de carbono implica un aumento de la dureza de la superficie (como se explica en la Sección 12.7) que, a su vez, conduce a un incremento de la resistencia al desgaste del engranaje [...]"

seccion 12.5.3; segundo párrafo: "[...] Casi siempre se utilizan en la condición templada y revenida, en la cual son especialmente resistentes al desgaste y capaces de adquirir la forma de herramienta de corte. [...]"

seccion 12.6.3; sextimo párrafo: "[...] Su aplicación se limita a componentes de gran dureza y resistencia al desgaste y sin ductilidad, como por ejemplo los cilindros de los trenes de laminación.[...]"

seccion 14.17.1; segundo párrafo: "[...] aumentando así el rendimiento del combustible utilizado; excelentes resistencias al desgaste y a la corrosión; menores pérdidas por fricción; [...]"

seccion 14.17.1-Página 38; segundo párrafo: "[...] La resistencia al desgaste y/o al deterioro 'a temperaturas elevadas de partes metálicas de los motores de combustión que se utilizan actualmente ha sido mejorada significativamente utilizando recubrimientos cerámicos.[...]"

seccion 23.7-Página 42; primer párrafo: "[...] puesto que las superficies articulares de la bola y de la cavidad acetabular rozan una con otra, el desgaste de estas superficies es minimizado empleando materiales muy duros.[...]"

seccion 23.8-Página 43; segundo párrafo: "[...] La cerámica seleccionada es el óxido de aluminio de alta pureza, el cual es más duro y más resistente al desgaste y genera tensiones de fricción inferiores en la articulación. [...]"

Glosario-pag 867; A: "[...] abrasivo. Material duro y resistente al desgaste (generalmente cerámico) que se utiliza para mecanizar, triturar o cortar otros materiales [...]"\\[0.1cm]

\textbf{Corrosión:}

Sección 13.1, Aleaciones Ferrosas-pag 433; primer párrafo: "[...] El principal inconveniente de muchas aleaciones férricas es la susceptibilidad a la corrosión [...]"

Sección 13.2, Resumen-pag 457; primer párrafo: "[...] otros elementos de aleación que los hacen susceptibles a tratamientos térmicos, a una mejora de las propiedades mecánicas y/o más resistentes a la corrosión [...]"

Sección 13.2, Introducción-pag 539; segundo párrafo: "[...] la ingeniería aeronáutica persigue materiales estructurales que tengan bajas densidades pero que sean rígidos y resistentes a la abrasión, a los impactos y a la corrosión [...]"

Cap 10-pag 585; ¿POR QUÉ ESTUDIAR Fabricación y procesado de materiales en ingeniería?; segundo párrafo: "[...]"Asimismo, algunos aceros inoxidables se hacen susceptibles a la corrosión intergranular "[...]"

Cap 17-pag 639; ¿POR QUÉ ESTUDIAR Corrosión y degradación de los materiales?; primer párrafo: "[...]"El conocimiento de los tipos de corrosión y de
degradación, y la comprensión de sus mecanismos"[...]"

Sección 18.1, Introducción-pag 639; Segundo párrafo: "[...] En losmetales hay pérdida de material por disolución (corrosión) o por formación de una capa o película no metálica (oxidación). [...]"\\[0.1cm]

\textbf{Grieta:}

Cap 10-pag 285; segundo párrafo: "[...] Esta grieta comenzó en algún tipo de pequeña entalla o hendidura aguda [...]"

Sección 10.2-pag 287; tercer párrafo: "[...] Cualquier proceso de fractura está compuesto de dos etapas –formación y propagación de la grieta– en respuesta a una tensión impuesta [...]"

Sección 10.2-pag 287; primer párrafo: "[...] Una grieta de este tipo normalmente se denomina estable; es decir, se resiste a extenderse a menos que haya un aumento de la tensión aplicada [...]"

Sección 10.2-pag 287; primer párrafo: "[...] Las grietas de este tipo se denominan inestables, y la propagación de la grieta, una vez iniciada, continúa espontáneamente sin que se requiera un aumento en la magnitud de la tensión aplicada [...]"

seccion 11.1; tercer párrafo: "[...] Si la velocidad de cambio de temperatura es grande, se genera un gradiente de temperatura que induce tensiones internas que pueden conducir a deformaciones e incluso al agrietamiento.[...]"

seccion 14.17.3-Página 40; cuarto párrafo: "[...] El encogimiento excesivo o demasiado rápido puede provocar agrietamiento y/o distorsión y por tanto una pieza sin	valor.[...]"

Sección 17.2-pag 586; primer párrafo: "[...] La mayoría de los materiales metálicos son especialmente susceptibles al hechurado, puesto que son moderadamente dúctiles y capaces de experimentar deformación permanente sin romperse ni agrietarse [...]"\\[0.1cm]



\end{document}